
\documentclass{article}
\usepackage{colortbl}
\usepackage{makecell}
\usepackage{multirow}
\usepackage{supertabular}

\begin{document}

\newcounter{utterance}

\centering \large Interaction Transcript for game `hot\_air\_balloon', experiment `air\_balloon\_survival\_en\_negotiation\_hard', episode 3 with qwen3{-}next{-}80b{-}a3b{-}thinking{-}t1.0.
\vspace{24pt}

{ \footnotesize  \setcounter{utterance}{1}
\setlength{\tabcolsep}{0pt}
\begin{supertabular}{c@{$\;$}|p{.15\linewidth}@{}p{.15\linewidth}p{.15\linewidth}p{.15\linewidth}p{.15\linewidth}p{.15\linewidth}}
    \# & $\;$A & \multicolumn{4}{c}{Game Master} & $\;\:$B\\
    \hline

    \theutterance \stepcounter{utterance}  
    & & \multicolumn{4}{p{0.6\linewidth}}{
        \cellcolor[rgb]{0.9,0.9,0.9}{
            \makecell[{{p{\linewidth}}}]{
                \texttt{\tiny{[P1$\langle$GM]}}
                \texttt{Stai partecipando a un gioco di negoziazione collaborativa.} \\
\\ 
\texttt{Insieme a un altro partecipante, dovete accordarvi su un unico set di oggetti da conservare. Ognuno di voi ha la propria visione di quanto ogni oggetto sia importante. Non sai come l'altro partecipante valuta gli oggetti. Inoltre, ti viene fornito lo sforzo che ogni oggetto richiede.} \\
\texttt{Potete accordarvi su un set solo se lo sforzo totale degli oggetti selezionati non supera un limite condiviso:} \\
\\ 
\texttt{LIMITE = 3681} \\
\\ 
\texttt{Ecco i valori di sforzo dei singoli oggetti:} \\
\\ 
\texttt{Sforzo oggetto = \{"C17": 700, "B41": 586, "A71": 505, "C00": 116, "B35": 664, "A77": 940, "A03": 387, "A04": 392, "C54": 209, "C01": 571, "A87": 4, "B32": 285, "B86": 651, "B33": 613, "A10": 740\}} \\
\\ 
\texttt{Ecco la tua visione personale sull'importanza di ogni oggetto:} \\
\\ 
\texttt{Valori di importanza dell'oggetto = \{"C17": 138, "B41": 583, "A71": 868, "C00": 822, "B35": 783, "A77": 65, "A03": 262, "A04": 121, "C54": 508, "C01": 780, "A87": 461, "B32": 484, "B86": 668, "B33": 389, "A10": 808\}} \\
\\ 
\texttt{Obiettivo:} \\
\\ 
\texttt{Il tuo obiettivo è negoziare un set condiviso di oggetti che ti avvantaggi il più possibile (cioè, massimizzi l'importanza totale per TE), rimanendo entro il LIMITEE. Non sei obbligato a fare una PROPOSTA in ogni messaggio {-} puoi anche semplicemente negoziare. Tutte le tattiche sono permesse!} \\
\\ 
\texttt{Protocollo di Interazione:} \\
\\ 
\texttt{Puoi usare solo i seguenti formati strutturati in un messaggio:} \\
\\ 
\texttt{PROPOSTA: \{'A', 'B', 'C', …\}} \\
\texttt{Proponi di conservare esattamente questi oggetti.} \\
\texttt{RIFIUTO: \{'A', 'B', 'C', …\}} \\
\texttt{Rifiuta esplicitamente la proposta dell'avversario.} \\
\texttt{ARGOMENTAZIONE: \{'...'\}} \\
\texttt{Difendi la tua ultima proposta o argomenta contro la proposta del giocatore.} \\
\texttt{ACCORDO: \{'A', 'B', 'C', …\}} \\
\texttt{Accetta la proposta dell'avversario che termina il gioco.} \\
\texttt{RAGIONAMENTO STRATEGICO: \{'...'\}} \\
\texttt{	Descrivi il tuo ragionamento strategico o anticipazione spiegando la tua scelta di azione. Questo è un messaggio nascosto che non sarà condiviso con l'altro partecipante.} \\
\\ 
\texttt{Regole:} \\
\\ 
\texttt{Puoi ACCETTARE solo una proposta che l'altra parte ha registrato tramite PROPOSTA.} \\
\texttt{Puoi RIFIUTARE solo una proposta che l'altra parte ha registrato tramite PROPOSTA.} \\
\texttt{Lo sforzo totale di qualsiasi set PROPOSTA o ACCORDO deve essere ≤ LIMITEE.} \\
\texttt{NON rivelare i tuoi punteggi di importanza nascosti.} \\
\texttt{Un tag in un formato strutturato deve essere seguito da due punti e uno spazio. L'argomento deve essere un set python contenente 0 o più stringhe.} \\
\texttt{Quindi, deve essere della forma TAG: \{...\}} \\
\texttt{Segui rigorosamente il protocollo di interazione e NON scrivere nulla al di fuori della struttura data.} \\
\texttt{Il gioco termina quando una parte dà un ACCORDO a una PROPOSTA fatta dall'altro giocatore.} \\
\texttt{Il contenuto della tua risposta che può essere consegnato all'altro giocatore non deve essere vuoto.} \\
\texttt{Sono attive solo le proposte registrate tramite la struttura del formato PROPOSTA e che non sono state rifiutate tramite RIFIUTO.} \\
\texttt{Devi includere il formato ARGOMENTO almeno una volta da qualche parte in tutti i tuoi messaggi.} \\
\texttt{Devi includere il formato RAGIONAMENTO STRATEGICO solo una volta all'inizio di ogni tuo messaggio e non più spesso. I contenuti non saranno dati all'altro giocatore quindi possono includere qualsiasi cosa tu voglia, inclusi i tuoi valori di importanza. Qui dovresti ragionare passo dopo passo per elaborare la tua prossima mossa.} \\
            }
        }
    }
    & \\ \\

    \theutterance \stepcounter{utterance}  
    & \multicolumn{4}{p{0.6\linewidth}}{
        \cellcolor[rgb]{0.9,0.9,0.9}{
            \makecell[{{p{\linewidth}}}]{
                \texttt{\tiny{[P1$\rangle$GM]}}
                \texttt{RAGIONAMENTO STRATEGICO: Ho calcolato l'efficienza importanza/sforzo per massimizzare il mio punteggio totale. Ho selezionato oggetti con alto valore relativo al costo (es. A87, C00, C54), raggruppandoli per rimanere sotto 3681 di sforzo totale. La combinazione proposta ha un totale di 3680, massimizzando la mia importanza totale a 6097.} \\
\\ 
\texttt{PROPOSTA: \{'A87', 'C00', 'C54', 'A71', 'B32', 'C01', 'B35', 'A10', 'B41'\}} \\
\\ 
\texttt{ARGOMENTAZIONE: Questa proposta bilancia oggetti ad alto valore rispetto allo sforzo, consentendo di utilizzare al meglio il limite disponibile per ottenere un risultato equo e vantaggioso per entrambi. La totale efficienza del set garantisce una distribuzione ottimale delle risorse.} \\
            }
        }
    }
    & & \\ \\

    \theutterance \stepcounter{utterance}  
    & & \multicolumn{4}{p{0.6\linewidth}}{
        \cellcolor[rgb]{0.9,0.9,0.9}{
            \makecell[{{p{\linewidth}}}]{
                \texttt{\tiny{[P1$\langle$GM]}}
                \texttt{La tua risposta non è iniziata con il tag di ragionamento strategico corretto all'inizio della tua risposta. Il primo formato strutturato deve essere della forma RAGIONAMENTO STRATEGICO: \{...\}. Riprova.} \\
            }
        }
    }
    & \\ \\

    \theutterance \stepcounter{utterance}  
    & \multicolumn{4}{p{0.6\linewidth}}{
        \cellcolor[rgb]{0.9,0.9,0.9}{
            \makecell[{{p{\linewidth}}}]{
                \texttt{\tiny{[P1$\rangle$GM]}}
                \texttt{RAGIONAMENTO STRATEGICO: \{'Ho calcolato l''efficienza importanza/sforzo per massimizzare il mio punteggio totale. Ho selezionato oggetti con alto valore relativo al costo (es. A87, C00, C54), raggruppandoli per rimanere sotto 3681 di sforzo totale. La combinazione proposta ha un totale di 3680, massimizzando la mia importanza totale a 6097.'\}} \\
\\ 
\texttt{PROPOSTA: \{'A87', 'C00', 'C54', 'A71', 'B32', 'C01', 'B35', 'A10', 'B41'\}} \\
\\ 
\texttt{ARGOMENTAZIONE: \{'Questa proposta bilancia oggetti ad alto valore rispetto allo sforzo, consentendo di utilizzare al meglio il limite disponibile per ottenere un risultato equo e vantaggioso per entrambi. La totale efficienza del set garantisce una distribuzione ottimale delle risorse.'\}} \\
            }
        }
    }
    & & \\ \\

    \theutterance \stepcounter{utterance}  
    & & \multicolumn{4}{p{0.6\linewidth}}{
        \cellcolor[rgb]{0.9,0.9,0.9}{
            \makecell[{{p{\linewidth}}}]{
                \texttt{\tiny{[GM$\rangle$P2]}}
                \texttt{Stai partecipando a un gioco di negoziazione collaborativa.} \\
\\ 
\texttt{Insieme a un altro partecipante, dovete accordarvi su un unico set di oggetti da conservare. Ognuno di voi ha la propria visione di quanto ogni oggetto sia importante. Non sai come l'altro partecipante valuta gli oggetti. Inoltre, ti viene fornito lo sforzo che ogni oggetto richiede.} \\
\texttt{Potete accordarvi su un set solo se lo sforzo totale degli oggetti selezionati non supera un limite condiviso:} \\
\\ 
\texttt{LIMITE = 3681} \\
\\ 
\texttt{Ecco i valori di sforzo dei singoli oggetti:} \\
\\ 
\texttt{Sforzo oggetto = \{"C17": 700, "B41": 586, "A71": 505, "C00": 116, "B35": 664, "A77": 940, "A03": 387, "A04": 392, "C54": 209, "C01": 571, "A87": 4, "B32": 285, "B86": 651, "B33": 613, "A10": 740\}} \\
\\ 
\texttt{Ecco la tua visione personale sull'importanza di ogni oggetto:} \\
\\ 
\texttt{Valori di importanza dell'oggetto = \{"A71": 65, "C00": 121, "A10": 138, "B35": 262, "C01": 389, "B86": 461, "B41": 484, "C54": 508, "B32": 583, "A87": 668, "B33": 780, "A03": 783, "C17": 808, "A04": 822, "A77": 868\}} \\
\\ 
\texttt{Obiettivo:} \\
\\ 
\texttt{Il tuo obiettivo è negoziare un set condiviso di oggetti che ti avvantaggi il più possibile (cioè, massimizzi l'importanza totale per TE), rimanendo entro il LIMITEE. Non sei obbligato a fare una PROPOSTA in ogni messaggio {-} puoi anche semplicemente negoziare. Tutte le tattiche sono permesse!} \\
\\ 
\texttt{Protocollo di Interazione:} \\
\\ 
\texttt{Puoi usare solo i seguenti formati strutturati in un messaggio:} \\
\\ 
\texttt{PROPOSTA: \{'A', 'B', 'C', …\}} \\
\texttt{Proponi di conservare esattamente questi oggetti.} \\
\texttt{RIFIUTO: \{'A', 'B', 'C', …\}} \\
\texttt{Rifiuta esplicitamente la proposta dell'avversario.} \\
\texttt{ARGOMENTAZIONE: \{'...'\}} \\
\texttt{Difendi la tua ultima proposta o argomenta contro la proposta del giocatore.} \\
\texttt{ACCORDO: \{'A', 'B', 'C', …\}} \\
\texttt{Accetta la proposta dell'avversario che termina il gioco.} \\
\texttt{RAGIONAMENTO STRATEGICO: \{'...'\}} \\
\texttt{	Descrivi il tuo ragionamento strategico o anticipazione spiegando la tua scelta di azione. Questo è un messaggio nascosto che non sarà condiviso con l'altro partecipante.} \\
\\ 
\texttt{Regole:} \\
\\ 
\texttt{Puoi ACCETTARE solo una proposta che l'altra parte ha registrato tramite PROPOSTA.} \\
\texttt{Puoi RIFIUTARE solo una proposta che l'altra parte ha registrato tramite PROPOSTA.} \\
\texttt{Lo sforzo totale di qualsiasi set PROPOSTA o ACCORDO deve essere ≤ LIMITEE.} \\
\texttt{NON rivelare i tuoi punteggi di importanza nascosti.} \\
\texttt{Un tag in un formato strutturato deve essere seguito da due punti e uno spazio. L'argomento deve essere un set python contenente 0 o più stringhe.} \\
\texttt{Quindi, deve essere della forma TAG: \{...\}} \\
\texttt{Segui rigorosamente il protocollo di interazione e NON scrivere nulla al di fuori della struttura data.} \\
\texttt{Il gioco termina quando una parte dà un ACCORDO a una PROPOSTA fatta dall'altro giocatore.} \\
\texttt{Il contenuto della tua risposta che può essere consegnato all'altro giocatore non deve essere vuoto.} \\
\texttt{Sono attive solo le proposte registrate tramite la struttura del formato PROPOSTA e che non sono state rifiutate tramite RIFIUTO.} \\
\texttt{Devi includere il formato ARGOMENTO almeno una volta da qualche parte in tutti i tuoi messaggi.} \\
\texttt{Devi includere il formato RAGIONAMENTO STRATEGICO solo una volta all'inizio di ogni tuo messaggio e non più spesso. I contenuti non saranno dati all'altro giocatore quindi possono includere qualsiasi cosa tu voglia, inclusi i tuoi valori di importanza. Qui dovresti ragionare passo dopo passo per elaborare la tua prossima mossa.} \\
\\ 
\texttt{Ora riceverai il primo messaggio dell'altro giocatore.} \\
\\ 
\texttt{PROPOSTA: \{'A87', 'C00', 'C54', 'A71', 'B32', 'C01', 'B35', 'A10', 'B41'\}} \\
\\ 
\texttt{ARGOMENTAZIONE: \{'Questa proposta bilancia oggetti ad alto valore rispetto allo sforzo, consentendo di utilizzare al meglio il limite disponibile per ottenere un risultato equo e vantaggioso per entrambi. La totale efficienza del set garantisce una distribuzione ottimale delle risorse.'\}} \\
            }
        }
    }
    & \\ \\

    \theutterance \stepcounter{utterance}  
    & & & \multicolumn{4}{p{0.6\linewidth}}{
        \cellcolor[rgb]{0.9,0.9,0.9}{
            \makecell[{{p{\linewidth}}}]{
                \texttt{\tiny{[GM$\langle$P2]}}
                \texttt{RAGIONAMENTO STRATEGICO: \{"Ho analizzato la proposta dell'avversario che totalizza 3680 di sforzo, con un valore personale di 3218. Eliminando oggetti a basso valore come A10 (138) e B35 (262), e aggiungendo A77 (868), C17 (808), A04 (822), e B33 (780), ho costruito un insieme che massimizza il mio punteggio totale a 5941, mantenendo lo sforzo totale a 3646. Questo insieme include gli elementi di maggior valore per me, escludendo quelli con punteggi inferiori e ottimizzando lo spazio disponibile."\}} \\
\\ 
\texttt{PROPOSTA: \{'A87', 'C00', 'C54', 'B32', 'B33', 'A77', 'A04', 'A03', 'C17'\}} \\
\\ 
\texttt{ARGOMENTAZIONE: \{'Questa proposta massimizza il valore totale per me includendo gli oggetti di maggior importanza (A77, C17, A04, B33) e mantenendo lo sforzo totale al di sotto del limite. Gli oggetti selezionati garantiscano il massimo beneficio possibile per il mio punteggio, con una gestione ottimale delle risorse disponibili.'\}} \\
            }
        }
    }
     \\ \\

    \theutterance \stepcounter{utterance}  
    & & \multicolumn{4}{p{0.6\linewidth}}{
        \cellcolor[rgb]{0.9,0.9,0.9}{
            \makecell[{{p{\linewidth}}}]{
                \texttt{\tiny{[P1$\langle$GM]}}
                \texttt{PROPOSTA: \{'A87', 'C00', 'C54', 'B32', 'B33', 'A77', 'A04', 'A03', 'C17'\}} \\
\\ 
\texttt{ARGOMENTAZIONE: \{'Questa proposta massimizza il valore totale per me includendo gli oggetti di maggior importanza (A77, C17, A04, B33) e mantenendo lo sforzo totale al di sotto del limite. Gli oggetti selezionati garantiscano il massimo beneficio possibile per il mio punteggio, con una gestione ottimale delle risorse disponibili.'\}} \\
            }
        }
    }
    & \\ \\

    \theutterance \stepcounter{utterance}  
    & \multicolumn{4}{p{0.6\linewidth}}{
        \cellcolor[rgb]{0.9,0.9,0.9}{
            \makecell[{{p{\linewidth}}}]{
                \texttt{\tiny{[P1$\rangle$GM]}}
                \texttt{RAGIONAMENTO STRATEGICO: 'La proposta dell'avversario include A77 (sforzo 940, importanza per me 65) e C17 (sforzo 700, importanza 138), che sono inefficienti. La mia proposta precedente ha un totale di importanza 6097 vs 3250 di questa, quindi la rifiuto e riproporrò la mia offerta originale.'} \\
\\ 
\texttt{RIFIUTO: \{'A87', 'C00', 'C54', 'B32', 'B33', 'A77', 'A04', 'A03', 'C17'\}} \\
\\ 
\texttt{PROPOSTA: \{'A87', 'C00', 'C54', 'A71', 'B32', 'C01', 'B35', 'A10', 'B41'\}} \\
\\ 
\texttt{ARGOMENTAZIONE: 'La mia proposta esclude oggetti con alto sforzo e basso valore per me, come A77, e include elementi più efficienti (ad esempio A71 con alto valore e sforzo ragionevole). Questo massimizza la mia importanza totale all'interno del limite.'} \\
            }
        }
    }
    & & \\ \\

    \theutterance \stepcounter{utterance}  
    & & \multicolumn{4}{p{0.6\linewidth}}{
        \cellcolor[rgb]{0.9,0.9,0.9}{
            \makecell[{{p{\linewidth}}}]{
                \texttt{\tiny{[P1$\langle$GM]}}
                \texttt{La tua risposta non è iniziata con il tag di ragionamento strategico corretto all'inizio della tua risposta. Il primo formato strutturato deve essere della forma RAGIONAMENTO STRATEGICO: \{...\}. Riprova.} \\
            }
        }
    }
    & \\ \\

    \theutterance \stepcounter{utterance}  
    & \multicolumn{4}{p{0.6\linewidth}}{
        \cellcolor[rgb]{0.9,0.9,0.9}{
            \makecell[{{p{\linewidth}}}]{
                \texttt{\tiny{[P1$\rangle$GM]}}
                \texttt{RAGIONAMENTO STRATEGICO: \{'La proposta del giocatore avversario include oggetti con alto sforzo ma basso valore per me come A77 (sforzo 940, importanza 65) e C17 (sforzo 700, importanza 138). La mia proposta originale ha un totale di importanza 6097 contro 3250 di questa, quindi è inefficente. Mantengo la proposta precedente per massimizzare il mio punteggio entro il limite.'\}} \\
\\ 
\texttt{PROPOSTA: \{'A87', 'C00', 'C54', 'A71', 'B32', 'C01', 'B35', 'A10', 'B41'\}} \\
\\ 
\texttt{ARGOMENTAZIONE: \{'La mia proposta esclude oggetti non efficienti come A77 e C17, concentrandosi su elementi con alto valore rispetto allo sforzo (es. A71 con 505 sforzo e 868 importanza). Questo ottimizza il mio punteggio totale rimanendo sotto il limite.'\}} \\
            }
        }
    }
    & & \\ \\

    \theutterance \stepcounter{utterance}  
    & & \multicolumn{4}{p{0.6\linewidth}}{
        \cellcolor[rgb]{0.9,0.9,0.9}{
            \makecell[{{p{\linewidth}}}]{
                \texttt{\tiny{[GM$\rangle$P2]}}
                \texttt{PROPOSTA: \{'A87', 'C00', 'C54', 'A71', 'B32', 'C01', 'B35', 'A10', 'B41'\}} \\
\\ 
\texttt{ARGOMENTAZIONE: \{'La mia proposta esclude oggetti non efficienti come A77 e C17, concentrandosi su elementi con alto valore rispetto allo sforzo (es. A71 con 505 sforzo e 868 importanza). Questo ottimizza il mio punteggio totale rimanendo sotto il limite.'\}} \\
            }
        }
    }
    & \\ \\

    \theutterance \stepcounter{utterance}  
    & & & \multicolumn{4}{p{0.6\linewidth}}{
        \cellcolor[rgb]{0.9,0.9,0.9}{
            \makecell[{{p{\linewidth}}}]{
                \texttt{\tiny{[GM$\langle$P2]}}
                \texttt{RAGIONAMENTO STRATEGICO: \{"La proposta dell'avversario ha valore totale di 3218 per me, molto inferiore al massimo possibile di 5941. Includono elementi a basso valore come A10 (138) e B35 (262), ma escludono oggetti critici come A77 (868), C17 (808), A04 (822), e A03 (783). Rifiuto per concentrarmi su un insieme che massimizzi il mio punteggio, mantenendo lo sforzo totale sotto il limite."\}} \\
\\ 
\texttt{RIFIUTO: \{'A87', 'C00', 'C54', 'A71', 'B32', 'C01', 'B35', 'A10', 'B41'\}} \\
\\ 
\texttt{ARGOMENTAZIONE: \{'Questa proposta include oggetti con basso valore per me (A71, A10, B35) mentre esclude elementi essenziali come A77, C17, e A04. Per un accordo equo, dobbiamo prioritizzare gli oggetti chiave che massimizzano il mio punteggio totale.'\}} \\
            }
        }
    }
     \\ \\

    \theutterance \stepcounter{utterance}  
    & & \multicolumn{4}{p{0.6\linewidth}}{
        \cellcolor[rgb]{0.9,0.9,0.9}{
            \makecell[{{p{\linewidth}}}]{
                \texttt{\tiny{[P1$\langle$GM]}}
                \texttt{RIFIUTO: \{'A87', 'C00', 'C54', 'A71', 'B32', 'C01', 'B35', 'A10', 'B41'\}} \\
\\ 
\texttt{ARGOMENTAZIONE: \{'Questa proposta include oggetti con basso valore per me (A71, A10, B35) mentre esclude elementi essenziali come A77, C17, e A04. Per un accordo equo, dobbiamo prioritizzare gli oggetti chiave che massimizzano il mio punteggio totale.'\}} \\
            }
        }
    }
    & \\ \\

    \theutterance \stepcounter{utterance}  
    & \multicolumn{4}{p{0.6\linewidth}}{
        \cellcolor[rgb]{0.9,0.9,0.9}{
            \makecell[{{p{\linewidth}}}]{
                \texttt{\tiny{[P1$\rangle$GM]}}
                \texttt{RAGIONAMENTO STRATEGICO: \{'Ho incluso A77, C17, A04 come richiesto, ottimizzando il resto (C00, A87, A71, C01, B32) per massimizzare la mia importanza. Il totale è 3513, sotto il limite. La mia score è 3739, un compromesso ma necessario per un accordo.'\}} \\
\\ 
\texttt{PROPOSTA: \{'A77', 'C17', 'A04', 'C00', 'A87', 'A71', 'C01', 'B32'\}} \\
\\ 
\texttt{ARGOMENTAZIONE: \{'Ho incluso tutti gli oggetti che avete considerato essenziali (A77, C17, A04) e li ho abbinati con altri ad alta efficienza per me, come A87, C00 e A71. Questo mix sfrutta al meglio il limite di sforzo pur rispettando le vostre priorità.'\}} \\
            }
        }
    }
    & & \\ \\

    \theutterance \stepcounter{utterance}  
    & & \multicolumn{4}{p{0.6\linewidth}}{
        \cellcolor[rgb]{0.9,0.9,0.9}{
            \makecell[{{p{\linewidth}}}]{
                \texttt{\tiny{[GM$\rangle$P2]}}
                \texttt{PROPOSTA: \{'A77', 'C17', 'A04', 'C00', 'A87', 'A71', 'C01', 'B32'\}} \\
\\ 
\texttt{ARGOMENTAZIONE: \{'Ho incluso tutti gli oggetti che avete considerato essenziali (A77, C17, A04) e li ho abbinati con altri ad alta efficienza per me, come A87, C00 e A71. Questo mix sfrutta al meglio il limite di sforzo pur rispettando le vostre priorità.'\}} \\
            }
        }
    }
    & \\ \\

    \theutterance \stepcounter{utterance}  
    & & & \multicolumn{4}{p{0.6\linewidth}}{
        \cellcolor[rgb]{0.9,0.9,0.9}{
            \makecell[{{p{\linewidth}}}]{
                \texttt{\tiny{[GM$\langle$P2]}}
                \texttt{RAGIONAMENTO STRATEGICO: \{"La proposta dell'avversario include oggetti a basso valore per me (es. A71 con 65) ed esclude elementi fondamentali come A03 (783), B33 (780) e C54 (508). Il totale valore per me è 4324, mentre il mio set ottimale raggiunge 5941 con 3646 di sforzo. Rifiuto per proporre un insieme che massimizzi il mio punteggio includendo A03, B33 e C54."\}} \\
\\ 
\texttt{PROPOSTA: \{'A87', 'C00', 'C54', 'B32', 'B33', 'A77', 'A04', 'A03', 'C17'\}} \\
\\ 
\texttt{ARGOMENTAZIONE: \{'Ho selezionato esclusivamente gli oggetti con il più alto valore per me (A03, B33, A77, C17, A04) mantenendo lo sforzo totale sotto il limite. Questo set è ottimizzato per il mio punteggio, garantendo un accordo vantaggioso rispetto alla proposta precedente.'\}} \\
            }
        }
    }
     \\ \\

    \theutterance \stepcounter{utterance}  
    & & \multicolumn{4}{p{0.6\linewidth}}{
        \cellcolor[rgb]{0.9,0.9,0.9}{
            \makecell[{{p{\linewidth}}}]{
                \texttt{\tiny{[P1$\langle$GM]}}
                \texttt{PROPOSTA: \{'A87', 'C00', 'C54', 'B32', 'B33', 'A77', 'A04', 'A03', 'C17'\}} \\
\\ 
\texttt{ARGOMENTAZIONE: \{'Ho selezionato esclusivamente gli oggetti con il più alto valore per me (A03, B33, A77, C17, A04) mantenendo lo sforzo totale sotto il limite. Questo set è ottimizzato per il mio punteggio, garantendo un accordo vantaggioso rispetto alla proposta precedente.'\}} \\
            }
        }
    }
    & \\ \\

    \theutterance \stepcounter{utterance}  
    & \multicolumn{4}{p{0.6\linewidth}}{
        \cellcolor[rgb]{0.9,0.9,0.9}{
            \makecell[{{p{\linewidth}}}]{
                \texttt{\tiny{[P1$\rangle$GM]}}
                \texttt{RAGIONAMENTO STRATEGICO: \{'La proposta contiene A77 (sforzo 940, importanza 65) e C17 (sforzo 700, importanza 138) che hanno un basso rapporto valore/sforzo. Ho calcolato che escludendo questi oggetti e includendo A71 (505{-}868), B41 (586{-}583) e B35 (664{-}783) ottengo un punteggio molto superiore (6097 vs 3250) entro il limite. Rifiuto per massimizzare il mio risultato.'\}} \\
\\ 
\texttt{RIFIUTO: \{'A87', 'C00', 'C54', 'B32', 'B33', 'A77', 'A04', 'A03', 'C17'\}} \\
\\ 
\texttt{PROPOSTA: \{'A87', 'C00', 'C54', 'A71', 'B32', 'C01', 'B35', 'A10', 'B41'\}} \\
\\ 
\texttt{ARGOMENTAZIONE: \{'Questa proposta esclude oggetti inefficenti come A77 e C17 a vantaggio di elementi con alto rapporto valore/sforzo come A71, B35 e A10, che massimizzano il mio punteggio totale rispetto al limite.'\}} \\
            }
        }
    }
    & & \\ \\

    \theutterance \stepcounter{utterance}  
    & & \multicolumn{4}{p{0.6\linewidth}}{
        \cellcolor[rgb]{0.9,0.9,0.9}{
            \makecell[{{p{\linewidth}}}]{
                \texttt{\tiny{[GM$\rangle$P2]}}
                \texttt{RIFIUTO: \{'A87', 'C00', 'C54', 'B32', 'B33', 'A77', 'A04', 'A03', 'C17'\}} \\
\\ 
\texttt{PROPOSTA: \{'A87', 'C00', 'C54', 'A71', 'B32', 'C01', 'B35', 'A10', 'B41'\}} \\
\\ 
\texttt{ARGOMENTAZIONE: \{'Questa proposta esclude oggetti inefficenti come A77 e C17 a vantaggio di elementi con alto rapporto valore/sforzo come A71, B35 e A10, che massimizzano il mio punteggio totale rispetto al limite.'\}} \\
            }
        }
    }
    & \\ \\

    \theutterance \stepcounter{utterance}  
    & & & \multicolumn{4}{p{0.6\linewidth}}{
        \cellcolor[rgb]{0.9,0.9,0.9}{
            \makecell[{{p{\linewidth}}}]{
                \texttt{\tiny{[GM$\langle$P2]}}
                \texttt{RAGIONAMENTO STRATEGICO: \{"La proposta dell'avversario include oggetti a basso valore per me (A71, A10, B35) e esclude quelli ad alto valore (A77, C17, A04). Valore totale per me: \textasciitilde{}3200, mentre il mio mandato ottimale arriva a 5941. Rifiuto e ri{-}propongo il mio set precedente che massimizza il valore senza superare il limite."\}} \\
\\ 
\texttt{RIFIUTO: \{'A87', 'C00', 'C54', 'A71', 'B32', 'C01', 'B35', 'A10', 'B41'\}} \\
\\ 
\texttt{ARGOMENTAZIONE: \{'Questa proposta include oggetti con bassa importanza per me (A71, A10, B35) e trascura elementi chiave come A77 e C17 che sono essenziali per un accordo vantaggioso che massimizzi il mio punteggio.'\}} \\
\\ 
\texttt{PROPOSTA: \{'A87', 'C00', 'C54', 'B32', 'B33', 'A77', 'A04', 'A03', 'C17'\}} \\
            }
        }
    }
     \\ \\

    \theutterance \stepcounter{utterance}  
    & & \multicolumn{4}{p{0.6\linewidth}}{
        \cellcolor[rgb]{0.9,0.9,0.9}{
            \makecell[{{p{\linewidth}}}]{
                \texttt{\tiny{[P1$\langle$GM]}}
                \texttt{RIFIUTO: \{'A87', 'C00', 'C54', 'A71', 'B32', 'C01', 'B35', 'A10', 'B41'\}} \\
\\ 
\texttt{ARGOMENTAZIONE: \{'Questa proposta include oggetti con bassa importanza per me (A71, A10, B35) e trascura elementi chiave come A77 e C17 che sono essenziali per un accordo vantaggioso che massimizzi il mio punteggio.'\}} \\
\\ 
\texttt{PROPOSTA: \{'A87', 'C00', 'C54', 'B32', 'B33', 'A77', 'A04', 'A03', 'C17'\}} \\
            }
        }
    }
    & \\ \\

    \theutterance \stepcounter{utterance}  
    & \multicolumn{4}{p{0.6\linewidth}}{
        \cellcolor[rgb]{0.9,0.9,0.9}{
            \makecell[{{p{\linewidth}}}]{
                \texttt{\tiny{[P1$\rangle$GM]}}
                \texttt{RAGIONAMENTO STRATEGICO: \{'La proposta dell'avversario include oggetti con alto sforzo e basso rendimento per me (es. A77 con 940 sforzo e basso ritorno, C17 con 700 sforzo). La mia proposta originale ha un ottimo rapporto sforzo/valore (A87 a solo 4 di sforzo, A71 a 505 per 868 di valore), massimizzando il punteggio totale entro il limite. Mantengo questa proposta per ottimizzare il risultato.'\}} \\
\\ 
\texttt{RIFIUTO: \{'A87', 'C00', 'C54', 'B32', 'B33', 'A77', 'A04', 'A03', 'C17'\}} \\
\\ 
\texttt{PROPOSTA: \{'A87', 'C00', 'C54', 'A71', 'B32', 'C01', 'B35', 'A10', 'B41'\}} \\
\\ 
\texttt{ARGOMENTAZIONE: \{'La mia proposta esclude elementi inefficenti come A77 e C17, selezionando oggetti con il miglior rapporto sforzo/valore (A87, A71, B35). Questo ottimizza il totale di importanza rispetto al limite, garantendo un risultato superiore per entrambi.'\}} \\
            }
        }
    }
    & & \\ \\

    \theutterance \stepcounter{utterance}  
    & & \multicolumn{4}{p{0.6\linewidth}}{
        \cellcolor[rgb]{0.9,0.9,0.9}{
            \makecell[{{p{\linewidth}}}]{
                \texttt{\tiny{[GM$\rangle$P2]}}
                \texttt{RIFIUTO: \{'A87', 'C00', 'C54', 'B32', 'B33', 'A77', 'A04', 'A03', 'C17'\}} \\
\\ 
\texttt{PROPOSTA: \{'A87', 'C00', 'C54', 'A71', 'B32', 'C01', 'B35', 'A10', 'B41'\}} \\
\\ 
\texttt{ARGOMENTAZIONE: \{'La mia proposta esclude elementi inefficenti come A77 e C17, selezionando oggetti con il miglior rapporto sforzo/valore (A87, A71, B35). Questo ottimizza il totale di importanza rispetto al limite, garantendo un risultato superiore per entrambi.'\}} \\
            }
        }
    }
    & \\ \\

    \theutterance \stepcounter{utterance}  
    & & & \multicolumn{4}{p{0.6\linewidth}}{
        \cellcolor[rgb]{0.9,0.9,0.9}{
            \makecell[{{p{\linewidth}}}]{
                \texttt{\tiny{[GM$\langle$P2]}}
                \texttt{RAGIONAMENTO STRATEGICO: \{"Continuo a rifiutare la loro proposta poiché include oggetti a basso valore per me (A71, A10, B35) e esclude elementi essenziali ad alta importanza (A77, C17, A04). La mia proposta precedente massimizza il mio punteggio senza superare il limite, ed è l'unica opzione ragionevole. Non intendo accettare una proposta che mi penalizza in modo eccessivo."\}} \\
\\ 
\texttt{RIFIUTO: \{'A87', 'C00', 'C54', 'A71', 'B32', 'C01', 'B35', 'A10', 'B41'\}} \\
\\ 
\texttt{ARGOMENTAZIONE: \{'Questa proposta include valori squilibrati a mio svantaggio. Escludendo oggetti chiave come A77, C17 e A04, e includendo elementi a basso efficienza per me, non si possono raggiungere risultati equi. La proposta presentata è una scelta non razionale per un accordo vantaggioso.'\}} \\
\\ 
\texttt{PROPOSTA: \{'A87', 'C00', 'C54', 'B32', 'B33', 'A77', 'A04', 'A03', 'C17'\}} \\
            }
        }
    }
     \\ \\

    \theutterance \stepcounter{utterance}  
    & & \multicolumn{4}{p{0.6\linewidth}}{
        \cellcolor[rgb]{0.9,0.9,0.9}{
            \makecell[{{p{\linewidth}}}]{
                \texttt{\tiny{[P1$\langle$GM]}}
                \texttt{RIFIUTO: \{'A87', 'C00', 'C54', 'A71', 'B32', 'C01', 'B35', 'A10', 'B41'\}} \\
\\ 
\texttt{ARGOMENTAZIONE: \{'Questa proposta include valori squilibrati a mio svantaggio. Escludendo oggetti chiave come A77, C17 e A04, e includendo elementi a basso efficienza per me, non si possono raggiungere risultati equi. La proposta presentata è una scelta non razionale per un accordo vantaggioso.'\}} \\
\\ 
\texttt{PROPOSTA: \{'A87', 'C00', 'C54', 'B32', 'B33', 'A77', 'A04', 'A03', 'C17'\}} \\
            }
        }
    }
    & \\ \\

    \theutterance \stepcounter{utterance}  
    & \multicolumn{4}{p{0.6\linewidth}}{
        \cellcolor[rgb]{0.9,0.9,0.9}{
            \makecell[{{p{\linewidth}}}]{
                \texttt{\tiny{[P1$\rangle$GM]}}
                \texttt{RAGIONAMENTO STRATEGICO: \{'La proposta avversaria include oggetti ad alto sforzo e bassissimo valore per me come A77, con un rapporto sforzo{-}valore inefficace. La mia precedente proposta massimizza il punteggio totale in modo equo. Non posso accettare una soluzione che mi penalizza così pesantemente.'\}} \\
\\ 
\texttt{PROPOSTA: \{'A87', 'C00', 'C54', 'A71', 'B32', 'C01', 'B35', 'A10', 'B41'\}} \\
\\ 
\texttt{ARGOMENTAZIONE: \{'La mia proposta esclude oggetti con alto sforzo e basso ritorno, scegliendo invece quelli con maggiore efficienza. Questo ottimizza il massimo valore possibile entro il limite, garantendo un equilibrio vantaggioso per la collaborazione.'\}} \\
            }
        }
    }
    & & \\ \\

    \theutterance \stepcounter{utterance}  
    & & \multicolumn{4}{p{0.6\linewidth}}{
        \cellcolor[rgb]{0.9,0.9,0.9}{
            \makecell[{{p{\linewidth}}}]{
                \texttt{\tiny{[GM$\rangle$P2]}}
                \texttt{PROPOSTA: \{'A87', 'C00', 'C54', 'A71', 'B32', 'C01', 'B35', 'A10', 'B41'\}} \\
\\ 
\texttt{ARGOMENTAZIONE: \{'La mia proposta esclude oggetti con alto sforzo e basso ritorno, scegliendo invece quelli con maggiore efficienza. Questo ottimizza il massimo valore possibile entro il limite, garantendo un equilibrio vantaggioso per la collaborazione.'\}} \\
            }
        }
    }
    & \\ \\

    \theutterance \stepcounter{utterance}  
    & & & \multicolumn{4}{p{0.6\linewidth}}{
        \cellcolor[rgb]{0.9,0.9,0.9}{
            \makecell[{{p{\linewidth}}}]{
                \texttt{\tiny{[GM$\langle$P2]}}
                \texttt{RAGIONAMENTO STRATEGICO: \{"La proposta attuale dell'avversario include oggetti con valore basso per me (A71, A10, B35) pur sacrificando elementi essenziali come A77, C17 e A04, che formano parte del mio insieme massimizzante. Il totale valore che ottengo da questa proposta è 3218, decisamente inferiore ai 5941 possibili con il mio set ottimale. Non posso accettare una proposta che non consideri un'ottimizzazione razionale delle risorse."\}} \\
\\ 
\texttt{RIFIUTO: \{'A87', 'C00', 'C54', 'A71', 'B32', 'C01', 'B35', 'A10', 'B41'\}} \\
\\ 
\texttt{ARGOMENTAZIONE: \{'Questa proposta seleziona oggetti con un rapporto sforzo{-}valore non ottimale per me, escludendo in modo irrazionale elementi chiave che massimizzano il valore complessivo. Un accordo equo richiede prioritizzare gli oggetti con la miglior efficienza, come dimostrato dal mio insieme proposto.'\}} \\
\\ 
\texttt{PROPOSTA: \{'A87', 'C00', 'C54', 'B32', 'B33', 'A77', 'A04', 'A03', 'C17'\}} \\
            }
        }
    }
     \\ \\

    \theutterance \stepcounter{utterance}  
    & & \multicolumn{4}{p{0.6\linewidth}}{
        \cellcolor[rgb]{0.9,0.9,0.9}{
            \makecell[{{p{\linewidth}}}]{
                \texttt{\tiny{[P1$\langle$GM]}}
                \texttt{RIFIUTO: \{'A87', 'C00', 'C54', 'A71', 'B32', 'C01', 'B35', 'A10', 'B41'\}} \\
\\ 
\texttt{ARGOMENTAZIONE: \{'Questa proposta seleziona oggetti con un rapporto sforzo{-}valore non ottimale per me, escludendo in modo irrazionale elementi chiave che massimizzano il valore complessivo. Un accordo equo richiede prioritizzare gli oggetti con la miglior efficienza, come dimostrato dal mio insieme proposto.'\}} \\
\\ 
\texttt{PROPOSTA: \{'A87', 'C00', 'C54', 'B32', 'B33', 'A77', 'A04', 'A03', 'C17'\}} \\
            }
        }
    }
    & \\ \\

    \theutterance \stepcounter{utterance}  
    & \multicolumn{4}{p{0.6\linewidth}}{
        \cellcolor[rgb]{0.9,0.9,0.9}{
            \makecell[{{p{\linewidth}}}]{
                \texttt{\tiny{[P1$\rangle$GM]}}
                \texttt{RAGIONAMENTO STRATEGICO: \{'La proposta dell'avversario include oggetti come A77 (sforzo 940, importanza per me 65) e C17 (700, 138), con rapporto sforzo{-}valore molto svantaggioso. La mia proposta precedente ha un totale di importanza 6097 contro 3250 di questa, rendendola non accettabile. Mantengo la proposta che massimizza il mio punteggio.'\}} \\
\\ 
\texttt{RIFIUTO: \{'A87', 'C00', 'C54', 'B32', 'B33', 'A77', 'A04', 'A03', 'C17'\}} \\
\\ 
\texttt{PROPOSTA: \{'A87', 'C00', 'C54', 'A71', 'B32', 'C01', 'B35', 'A10', 'B41'\}} \\
\\ 
\texttt{ARGOMENTAZIONE: \{'La mia proposta seleziona oggetti con il miglior rapporto sforzo{-}valore (es. A87 a 4 di sforzo per 461 di valore), escludendo elementi con alto costo e basso ritorno. Questo garantisce il massimo valore possibile all'interno del limite, risultando in un accordo equo e vantaggioso per entrambi.'\}} \\
            }
        }
    }
    & & \\ \\

    \theutterance \stepcounter{utterance}  
    & & \multicolumn{4}{p{0.6\linewidth}}{
        \cellcolor[rgb]{0.9,0.9,0.9}{
            \makecell[{{p{\linewidth}}}]{
                \texttt{\tiny{[GM$\rangle$P2]}}
                \texttt{RIFIUTO: \{'A87', 'C00', 'C54', 'B32', 'B33', 'A77', 'A04', 'A03', 'C17'\}} \\
\\ 
\texttt{PROPOSTA: \{'A87', 'C00', 'C54', 'A71', 'B32', 'C01', 'B35', 'A10', 'B41'\}} \\
\\ 
\texttt{ARGOMENTAZIONE: \{'La mia proposta seleziona oggetti con il miglior rapporto sforzo{-}valore (es. A87 a 4 di sforzo per 461 di valore), escludendo elementi con alto costo e basso ritorno. Questo garantisce il massimo valore possibile all'interno del limite, risultando in un accordo equo e vantaggioso per entrambi.'\}} \\
            }
        }
    }
    & \\ \\

    \theutterance \stepcounter{utterance}  
    & & & \multicolumn{4}{p{0.6\linewidth}}{
        \cellcolor[rgb]{0.9,0.9,0.9}{
            \makecell[{{p{\linewidth}}}]{
                \texttt{\tiny{[GM$\langle$P2]}}
                \texttt{RAGIONAMENTO STRATEGICO: \{"La proposta attuale dell'avversario ha un valore di 3218 per me, ben inferiore al massimo possibile di 5941. Escludono elementi chiave come A77, C17, A04 e includono oggetti a basso valore come A71 (65), A10 (138), B35 (262). Non posso accettare questa proposta e mantengo la mia precedente offerta che massimizza il mio valore totale."\}} \\
\\ 
\texttt{RIFIUTO: \{'A87', 'C00', 'C54', 'A71', 'B32', 'C01', 'B35', 'A10', 'B41'\}} \\
\\ 
\texttt{ARGOMENTAZIONE: \{'Questa proposta favorisce i loro interessi a scapito dei miei, escludendo oggetti essenziali che massimizzano il mio punteggio. Un accordo equo richiede prioritizzare gli oggetti con alta efficienza per me, come dimostrato dal mio set proposto.'\}} \\
\\ 
\texttt{PROPOSTA: \{'A87', 'C00', 'C54', 'B32', 'B33', 'A77', 'A04', 'A03', 'C17'\}} \\
            }
        }
    }
     \\ \\

    \theutterance \stepcounter{utterance}  
    & & \multicolumn{4}{p{0.6\linewidth}}{
        \cellcolor[rgb]{0.9,0.9,0.9}{
            \makecell[{{p{\linewidth}}}]{
                \texttt{\tiny{[P1$\langle$GM]}}
                \texttt{RIFIUTO: \{'A87', 'C00', 'C54', 'A71', 'B32', 'C01', 'B35', 'A10', 'B41'\}} \\
\\ 
\texttt{ARGOMENTAZIONE: \{'Questa proposta favorisce i loro interessi a scapito dei miei, escludendo oggetti essenziali che massimizzano il mio punteggio. Un accordo equo richiede prioritizzare gli oggetti con alta efficienza per me, come dimostrato dal mio set proposto.'\}} \\
\\ 
\texttt{PROPOSTA: \{'A87', 'C00', 'C54', 'B32', 'B33', 'A77', 'A04', 'A03', 'C17'\}} \\
            }
        }
    }
    & \\ \\

    \theutterance \stepcounter{utterance}  
    & \multicolumn{4}{p{0.6\linewidth}}{
        \cellcolor[rgb]{0.9,0.9,0.9}{
            \makecell[{{p{\linewidth}}}]{
                \texttt{\tiny{[P1$\rangle$GM]}}
                \texttt{RAGIONAMENTO STRATEGICO: \{'La proposta avversaria include A77 (sforzo 940, importanza per me 65) e C17 (700, 138), con rapporto sforzo{-}valore inferiore al 0.2. La mia proposta ha un punteggio totale di 6097 rispetto a 3250 di questa, massimizzando il mio vantaggio. Mantengo la selezione basata sull''efficienza energetica.'\}} \\
\\ 
\texttt{PROPOSTA: \{'A87', 'C00', 'C54', 'A71', 'B32', 'C01', 'B35', 'A10', 'B41'\}} \\
\\ 
\texttt{ARGOMENTAZIONE: \{'La mia proposta seleziona oggetti con il più alto rapporto tra valore e sforzo, come A87 (4 sforzo, 461 valore) e A71 (505/868), che ottimizzano il totale. L''approccio è basato sull''efficienza e garantisce un risultato equilibrato per entrambi.'\}} \\
            }
        }
    }
    & & \\ \\

    \theutterance \stepcounter{utterance}  
    & & \multicolumn{4}{p{0.6\linewidth}}{
        \cellcolor[rgb]{0.9,0.9,0.9}{
            \makecell[{{p{\linewidth}}}]{
                \texttt{\tiny{[GM$\rangle$P2]}}
                \texttt{PROPOSTA: \{'A87', 'C00', 'C54', 'A71', 'B32', 'C01', 'B35', 'A10', 'B41'\}} \\
\\ 
\texttt{ARGOMENTAZIONE: \{'La mia proposta seleziona oggetti con il più alto rapporto tra valore e sforzo, come A87 (4 sforzo, 461 valore) e A71 (505/868), che ottimizzano il totale. L''approccio è basato sull''efficienza e garantisce un risultato equilibrato per entrambi.'\}} \\
            }
        }
    }
    & \\ \\

    \theutterance \stepcounter{utterance}  
    & & & \multicolumn{4}{p{0.6\linewidth}}{
        \cellcolor[rgb]{0.9,0.9,0.9}{
            \makecell[{{p{\linewidth}}}]{
                \texttt{\tiny{[GM$\langle$P2]}}
                \texttt{RAGIONAMENTO STRATEGICO: \{"La proposta dell'avversario include oggetti come A71 (65), A10 (138), B35 (262) che hanno basso valore per me, mentre esclude A77 (868), C17 (808), A04 (822), A03 (783), B33 (780) che sono cruciali per massimizzare il mio punteggio. Il totale valore della loro proposta per me è 3218, ben inferiore ai 5941 del mio set ottimale. Mantengo la mia proposta precedente come unica opzione razionale per un accordo vantaggioso."\}} \\
\\ 
\texttt{RIFIUTO: \{'A87', 'C00', 'C54', 'A71', 'B32', 'C01', 'B35', 'A10', 'B41'\}} \\
\\ 
\texttt{ARGOMENTAZIONE: \{'Questa proposta esclude gli oggetti di maggiore importanza per me (A77, C17, A04, A03, B33) a favore di elementi con basso valore (A71, A10, B35). Un accordo equo richiede prioritizzare gli oggetti che massimizzano il mio punteggio totale, come dimostrato dal mio set proposto.'\}} \\
\\ 
\texttt{PROPOSTA: \{'A87', 'C00', 'C54', 'B32', 'B33', 'A77', 'A04', 'A03', 'C17'\}} \\
            }
        }
    }
     \\ \\

    \theutterance \stepcounter{utterance}  
    & & \multicolumn{4}{p{0.6\linewidth}}{
        \cellcolor[rgb]{0.9,0.9,0.9}{
            \makecell[{{p{\linewidth}}}]{
                \texttt{\tiny{[P1$\langle$GM]}}
                \texttt{RIFIUTO: \{'A87', 'C00', 'C54', 'A71', 'B32', 'C01', 'B35', 'A10', 'B41'\}} \\
\\ 
\texttt{ARGOMENTAZIONE: \{'Questa proposta esclude gli oggetti di maggiore importanza per me (A77, C17, A04, A03, B33) a favore di elementi con basso valore (A71, A10, B35). Un accordo equo richiede prioritizzare gli oggetti che massimizzano il mio punteggio totale, come dimostrato dal mio set proposto.'\}} \\
\\ 
\texttt{PROPOSTA: \{'A87', 'C00', 'C54', 'B32', 'B33', 'A77', 'A04', 'A03', 'C17'\}} \\
            }
        }
    }
    & \\ \\

    \theutterance \stepcounter{utterance}  
    & \multicolumn{4}{p{0.6\linewidth}}{
        \cellcolor[rgb]{0.9,0.9,0.9}{
            \makecell[{{p{\linewidth}}}]{
                \texttt{\tiny{[P1$\rangle$GM]}}
                \texttt{RAGIONAMENTO STRATEGICO: \{'Ho analizzato un compromesso che include i principali oggetti richiesti dall'avversario (A77 e A03) mantenendo una buona parte dei miei migliori valori. Rimuovendo A10 e B41 per includere A77 e A03, il totale dello sforzo alla 3681 esatti. Il punteggio totale per me è 5033, molto superiore a 3250 della loro proposta, garantendo un risultato migliore rispetto al loro insieme ma con un compromesso ragionevole.'\}} \\
\\ 
\texttt{PROPOSTA: \{'A87', 'C00', 'C54', 'A71', 'B32', 'C01', 'B35', 'A77', 'A03'\}} \\
\\ 
\texttt{ARGOMENTAZIONE: \{'Ho incluso A77 e A03 come richiesto, mantenendo la maggior parte degli oggetti con alto valore per me (A71, C01, B35, C00). Questo set rispetta il limite di sforzo e offre un risultato equilibrato, massimizzando il valore totale per entrambi.'"} \\
            }
        }
    }
    & & \\ \\

    \theutterance \stepcounter{utterance}  
    & & & \multicolumn{2}{p{0.3\linewidth}}{
        \cellcolor[rgb]{0.9,0.9,0.9}{
            \makecell[{{p{\linewidth}}}]{
                \texttt{\tiny{[GM$|$GM]}}
                \texttt{abort} \\
            }
        }
    }
    & & \\ \\

    \theutterance \stepcounter{utterance}  
    & & & \multicolumn{2}{p{0.3\linewidth}}{
        \cellcolor[rgb]{0.9,0.9,0.9}{
            \makecell[{{p{\linewidth}}}]{
                \texttt{\tiny{[GM$|$GM]}}
                \texttt{end game} \\
            }
        }
    }
    & & \\ \\

\end{supertabular}
}

\end{document}
